\documentclass[tikz,border=6pt]{standalone}
\usepackage{ctex}
\usepackage{amsmath}
\usetikzlibrary{shapes.geometric, positioning}

\begin{document}
\begin{tikzpicture}[
  every node/.style={font=\small},
  label node/.style={font=\footnotesize, text=white, align=center},
]

% ── 颜色定义 ──────────────────────────────────────────────────────────────────
% 从外到内：SDP > SOCP > QP > LP
\definecolor{colSDP}{RGB}{70,130,180}    % steel blue
\definecolor{colSOCP}{RGB}{60,179,113}   % medium sea green
\definecolor{colQP}{RGB}{255,165,0}      % orange
\definecolor{colLP}{RGB}{220,50,50}      % red

% ── 椭圆（从外到内） ──────────────────────────────────────────────────────────
% SDP（最外层）
\fill[colSDP!25] (0,0) ellipse (6.2cm and 4.0cm);
\draw[colSDP, thick] (0,0) ellipse (6.2cm and 4.0cm);

% SOCP
\fill[colSOCP!30] (0,-0.1) ellipse (4.6cm and 3.0cm);
\draw[colSOCP, thick] (0,-0.1) ellipse (4.6cm and 3.0cm);

% QP
\fill[colQP!30] (0,-0.2) ellipse (3.1cm and 2.1cm);
\draw[colQP, thick] (0,-0.2) ellipse (3.1cm and 2.1cm);

% LP（最内层）
\fill[colLP!30] (0,-0.25) ellipse (1.7cm and 1.2cm);
\draw[colLP, thick] (0,-0.25) ellipse (1.7cm and 1.2cm);

% ── 层标签（左上角） ─────────────────────────────────────────────────────────
\node[colSDP!80!black, font=\normalsize\bfseries] at (-4.8, 3.2) {SDP};
\node[colSOCP!70!black, font=\normalsize\bfseries] at (-3.4, 2.4) {SOCP};
\node[colQP!70!black, font=\normalsize\bfseries] at (-1.6, 1.55) {QP};
\node[colLP!70!black, font=\small\bfseries] at (0, -0.25) {LP};

% ── 典型例子标注 ──────────────────────────────────────────────────────────────
% LP 内部例子（紧凑）
\node[font=\tiny, align=center, text=colLP!70!black] at (0, -0.5) {
  L1回归\\
  Chebyshev逼近
};

% QP 环形区域（LP以外）
\node[font=\tiny, align=center] at (2.0, -0.3) {
  SVM\\（硬/软间隔）
};
\node[font=\tiny, align=center] at (-2.0, 0.1) {
  Markowitz\\投资组合
};

% SOCP 环形区域
\node[font=\tiny, align=center] at (3.8, -0.3) {
  鲁棒优化
};
\node[font=\tiny, align=center] at (-3.6, 0.3) {
  椭球约束
};

% SDP 环形区域
\node[font=\tiny, align=center] at (5.2, -0.3) {
  矩阵补全
};
\node[font=\tiny, align=center] at (-5.1, 0.3) {
  谱范数\\正则化
};
\node[font=\tiny, align=center] at (0, 3.4) {
  SDP松弛（非凸问题的凸近似）
};

% ── 包含关系说明（底部）────────────────────────────────────────────────────────
\node[font=\small, align=center] at (0, -3.7) {
  $\text{LP} \subset \text{QP} \subset \text{SOCP} \subset \text{SDP}$\\[2pt]
  {\footnotesize 均可在多项式时间内求解（内点法）}
};

\end{tikzpicture}
\end{document}
